\label{chap:83-06}

\section{Medida proyectiva y conservación evolutiva de la unidad}
\label{p12-83-c6-s1:chap:medida-proyectiva}

La topología de cilindros permite hablar de proximidad; una teoría de caminos
necesita además pesos. La conservación evolutiva de la unidad adquiere aquí
una forma probabilística exacta: cada cilindro reparte su peso entre sus
extensiones y la suma total permanece igual a uno. La unidad no queda fija en
una celda; se transporta a través del refinamiento.

\subsection{Pesos consistentes sobre niveles finitos}

Sea \(\mu_n\) una probabilidad sobre el conjunto finito
\(\mathcal S_n\). La familia es \emph{proyectivamente consistente} cuando
\begin{equation}
 \mu_n(a)
 =\sum_{\substack{b\in\mathcal S_{n+1}\\
                   \rho_{n+1,n}(b)=a}}
   \mu_{n+1}(b)
 \qquad(a\in\mathcal S_n).
 \label{p12-83-c6-s1:eq:consistencia-pesos}
\end{equation}
Esta igualdad expresa la conservación local de la unidad: el peso de una
historia truncada es la suma de los pesos de todos sus hijos.

\begin{theorem}[Medida genealógica sobre cilindros]
\label{p12-83-c6-s1:thm:medida-genealogica}
Toda familia de probabilidades finitas que satisface
\eqref{p12-83-c6-s1:eq:consistencia-pesos} determina una única probabilidad de Borel
\(\mu_\infty\) sobre \(\Omega_\infty\) tal que
\begin{equation}
 \mu_\infty([a]_n)=\mu_n(a).
 \label{p12-83-c6-s1:eq:medida-cilindro}
\end{equation}
\end{theorem}

\begin{proof}
\[
 \mathscr A
 :=\{\text{uniones finitas de cilindros}\}
\]
es un álgebra: la intersección de dos cilindros es vacía o un cilindro del
nivel más profundo. Todo elemento de \(\mathscr A\) puede refinarse a una
unión disjunta de cilindros de un mismo nivel. Si
\(A=\bigsqcup_{a\in I}[a]_n\), ponemos
\[
 \mu_0(A)=\sum_{a\in I}\mu_n(a).
\]
Iterar \eqref{p12-83-c6-s1:eq:consistencia-pesos} demuestra que el valor no depende del
nivel común elegido.

Si un cilindro \([a]_n\) es vacío, la ramificación finita y el lema de König
\cite{Koenig1936}
implican que existe una profundidad sin descendientes de \(a\). La
consistencia iterada da entonces \(\mu_n(a)=0\). Por ello \(\mu_0\) está bien
definida incluso para representaciones que contienen cilindros vacíos y es
finitamente aditiva.

Los cilindros son abiertos y cerrados. Si una sucesión decreciente de
elementos de \(\mathscr A\) tuviera intersección vacía, la compacidad
implicaría que alguno de ellos ya es vacío; por ello \(\mu_0\) es continua en
el vacío y constituye una premedida numerablemente aditiva. El teorema de extensión de
Carathéodory ---equivalentemente, la extensión de Kolmogórov para este sistema
proyectivo finito--- la prolonga a la \(\sigma\)-álgebra generada por los
cilindros \cite{Caratheodory1914,Kolmogorov1933}. Ésta es la Boreliana de
\(\Omega_\infty\), y la finitud de la
medida garantiza unicidad.
\end{proof}

La medida puede proceder de transiciones locales. Si
\(P_n(b\mid a)\ge0\) satisface
\[
 \sum_{\rho_{n+1,n}(b)=a}P_n(b\mid a)=1,
\]
entonces
\(\mu_{n+1}(b)=\mu_n(a)P_n(b\mid a)\) es consistente. La uniformidad es una
elección posible, no una obligación. Los pesos pueden depender de hoja,
acarreo, orientación, un coste de acción no negativo o un registro siempre
que preserven la suma. Por ejemplo,
\[
 P_n(b\mid a)
 =\frac{e^{-\beta a_n(b\mid a)}}
        {\sum_{b'\mapsto a}e^{-\beta a_n(b'\mid a)}},
 \qquad a_n\geq0,
\]
define pesos positivos normalizados. La fase ortogonal de acción pertenece al
kernel del \cref{chap:integral-genealogica}; no se inserta en estas
probabilidades positivas.

\subsection{Desintegración de la medida genealógica sobre fibras de igual valor arquimediano}

La evaluación arquimediana transporta la medida a \(K\):
\begin{equation}
 \nu:=\Val_*\mu_\infty,
 \qquad
 \nu(B)=\mu_\infty(\Val^{-1}(B)).
 \label{p12-83-c6-s1:eq:pushforward-arquimediano}
\end{equation}
La probabilidad visible de una región suma el peso de todas las genealogías que
la publican. Cuando dos secciones poseen el mismo valor, sus contribuciones se
agregan en \(\nu\), aunque permanezcan separadas en \(\mu_\infty\).

\begin{proposition}[Desintegración por valores]
Si \(K\) es un espacio estándar de Borel, existe una familia de medidas
condicionadas \(\{\mu_y\}_{y\in K}\), definida para \(\nu\)-casi todo \(y\),
tal que \(\mu_y\) está soportada en \(\mathfrak G_y=\Val^{-1}(y)\) y
\[
 \int_{\Omega_\infty}F\dd\mu_\infty
 =\int_K\left(\int_{\mathfrak G_y}F\dd\mu_y\right)\dd\nu(y)
\]
para toda función integrable \(F\).
\end{proposition}

Esta fórmula separa valor y genealogía en el nivel de la medida. La física que
dependa sólo del valor integra primero toda la fibra; una observable de memoria
puede distinguir distribuciones \(\mu_y\) diferentes sobre la misma
coordenada real.

\subsection{Condicionamiento como poda de futuros}

Sea \(A\subseteq\Omega_\infty\) un suceso observable con
\(\mu_\infty(A)>0\). La medida condicionada es
\begin{equation}
 \mu_\infty(B\mid A)
 =\frac{\mu_\infty(B\cap A)}{\mu_\infty(A)}.
 \label{p12-83-c6-s1:eq:condicionamiento-genealogico}
\end{equation}
Cuando \(A\) es la lectura de un registro, su intersección con cada cilindro
elimina las extensiones incompatibles. La normalización no crea nuevas ramas:
redistribuye la unidad entre las que han sobrevivido.

\begin{proposition}[Actualización compatible]
Si \(A\) es unión de cilindros a nivel \(m\), entonces para todo \(n\ge m\)
la distribución condicionada sobre \(\mathcal S_n\) se obtiene suprimiendo los
estados cuyos truncamientos no pertenecen a \(A\) y renormalizando los pesos
restantes. Las distribuciones condicionadas siguen satisfaciendo
\eqref{p12-83-c6-s1:eq:consistencia-pesos}.
\end{proposition}

\begin{proof}
La identidad es la regla de Bayes sobre conjuntos finitos. La suma de los
pesos de los hijos condicionados coincide con el peso condicionado del padre,
porque intersección y unión disjunta conmutan.
\end{proof}

\subsection{Conservación combinatoria y completación dimensional}

La misma ley aparece en la completación binomial:
\[
 1=\sum_{r=0}^{d}\binom dr
 \left(\frac9{10}\right)^{d-r}
 \left(\frac1{10}\right)^r.
\]
En dimensión tres,
\[
 1000=729+243+27+1,
\]
de modo que interior, caras nuevas, aristas extremas y vértice final forman
una partición de la unidad decimal. El paso
\(729+270=999\to729+271=1000\) distingue el estado anterior al cierre de la
unidad completada. En el sistema de caminos, la identidad análoga es
\[
 \mu([a]_n)=\sum_{b\mapsto a}\mu([b]_{n+1}).
\]
La unidad evoluciona al distribuirse por los estratos y se recupera al sumar
la genealogía completa.

\begin{figure}[htbp]
\centering
\begin{tikzpicture}[>=Latex,node distance=12mm and 17mm,
  box/.style={draw=hmtblue,rounded corners,fill=hmtlightblue,align=center,
              minimum width=28mm,minimum height=8mm},
  leaf/.style={draw=hmtgreen,rounded corners,fill=hmtlightgreen,align=center,
               minimum width=22mm,minimum height=7mm}]
\node[box] (a) {$[a]_n$\\peso $\mu_n(a)$};
\node[leaf,below left=of a] (b1) {$[b_1]_{n+1}$\\$\mu_{n+1}(b_1)$};
\node[leaf,below=of a] (b2) {$[b_2]_{n+1}$\\$\mu_{n+1}(b_2)$};
\node[leaf,below right=of a] (b3) {$[b_3]_{n+1}$\\$\mu_{n+1}(b_3)$};
\draw[->] (a)--(b1); \draw[->] (a)--(b2); \draw[->] (a)--(b3);
\node[below=20mm of b2,align=center,text width=.78\textwidth] (eq)
 {$\mu_n(a)=\mu_{n+1}(b_1)+\mu_{n+1}(b_2)+\mu_{n+1}(b_3)$:\\
  la unidad se conserva al refinar la genealogía.};
\end{tikzpicture}
\caption{Conservación proyectiva del peso sobre cilindros.}
\label{p12-83-c6-s1:fig:conservacion-peso-cilindros}
\end{figure}

\begin{statusbox}
La medida proyectiva se construye sobre el sistema inverso HMT.
El teorema de extensión es matemática clásica; la identificación de sus
cilindros con estados TPK es exacta una vez fijados los pesos de transición.
La selección de una medida física concreta pertenece al modelo dinámico y no
se deduce únicamente del censo de ramas.
\end{statusbox}

La medida proyectiva queda caracterizada por positividad, conservación del
peso entre padres e hijos y condicionamiento dentro del suceso observado. Una
realización experimental completa añade la preparación y el protocolo de
muestreo; con esos datos, la misma medida produce probabilidades de lectura
comparables con frecuencias.
